\section*{ВВЕДЕНИЕ}
\addcontentsline{toc}{section}{ВВЕДЕНИЕ}
Назначение археологии можно было бы объяснить простым любопытством, 
которое толкает человека к познанию и открытиям. Не последнюю роль играет 
стремление человека к изучению своих истоков, своей культуры. Человек без прошлого --- это человек без будущего. Совершенно определенно можно утверждать, что изучая прошлое, мы прогнозируем свое будущее. 

Археология невозможна без проведения полевых исследований, в основной массе 
это археологические разведки и раскопки, где все работы проводятся вручную с 
использованием огромного арсенала спецсредств. Опыт раскопок обязателен для любого археолога, даже применяющего иные методы исследований, которые, в свою очередь, обретают свою значимость только благодаря раскопкам. Только раскопками можно проверить идеи, модели и теории.

Важность раскопок на современном этапе развития науки и общества можно 
проиллюстрировать минимум тремя примерами. Первый пример касается времени 
появления современных людей и их сосуществования с неандертальцами. Второй пример относится к объектам XX века – исследованию так называемых <<ландшафтов террора>> --- т.е. концентрационных лагерей или массовых захоронений. Они позволяют восстановить детали, которые не отмечают письменные источники, как официальные документы, так и личные воспоминания. Третий пример не относится к научной археологии, но показывает наличие интереса общества к раскопкам. В частности, появления такого явления как геокэшинг, предполагающего поиск обычными людьми, <<неархеологами>>, специально спрятанных <<сокровищ>> в определенных местах, имеющих историческую значимость. 
Этот пример свидетельствует о том, что общество нуждается в возможности участвовать в раскопках и получать о них информацию.

Приведенные примеры показывают, что тенденции, как в самой археологии, так и в обществе связаны с возобновлением интереса к эмпирическому познанию, центральная роль в развитии которого, безусловно, должна принадлежать раскопкам.

В связи с ростом количества проводимых археологических экспедиций возникает 
необходимость создания информационной системы, позволяющей упорядочить и 
минимизировать время сбора необходимых данных для ее организации.

Цель работы --- разработка базы данных приложения по организации археологических экспедиций.

Для достижения поставленной цели необходимо решить следующие задачи:
\begin{itemize}[left=0.49cm, label*=---]
    \item формализовать задачу и данные;
    \item спроектировать базу данных;
    \item выбрать систему управления базами данных;
    \item реализовать спроектированную базу данных и необходимый интерфейс для взаимодействия с ней;
    \item исследовать зависимость времени выполнения запроса от наличия индексов.
\end{itemize}
